\documentclass[a4paper,11pt]{ltjsarticle}
\newcommand{\bookletname}{採点基準}
% LuaLaTeX / UTF-8
\usepackage{amsmath,amssymb,bm}
\usepackage[margin=23mm,top=22mm,bottom=22mm]{geometry}
\usepackage{enumitem,booktabs,array}
\usepackage{fancyhdr}
\usepackage[hidelinks,unicode]{hyperref}
\setlength{\parindent}{1\zw}
\setlength{\parskip}{0.35em}
\setlength{\headheight}{17pt}
\pagestyle{fancy}
\fancyhf{}
\fancyhead[L]{\small 第1回模試\quad 統計数理}
\fancyhead[R]{\small\bookletname}
\fancyfoot[C]{\small\thepage}
\renewcommand{\headrulewidth}{0.35pt}
\renewcommand{\footrulewidth}{0pt}
\setlist[enumerate,1]{label=〔\arabic*〕,leftmargin=3.5em,labelsep=0.5em,itemsep=1.1em,topsep=0.8em}
\newcommand{\E}{\operatorname{E}}
\newcommand{\Var}{\operatorname{Var}}
\newcommand{\Cov}{\operatorname{Cov}}
\newcommand{\Prob}{\operatorname{P}}
\newcommand{\ind}{\mathbf{1}}
\newcommand{\qhead}[1]{\section*{問#1}}
\newcommand{\ans}[1]{\subsection*{〔#1〕}}
\newcommand{\points}[1]{\hfill\textnormal{（#1点）}}
\newcommand{\continued}[1]{\noindent\textbf{問#1（続き）}\par\medskip}
\hypersetup{pdfauthor={Codex},pdfsubject={統計検定1級対策・自主学習用模擬試験}}

\hypersetup{pdftitle={第1回模試 統計数理 採点基準}}
\begin{document}
\begin{center}
{\LARGE\bfseries 第1回模試\quad 統計数理}\par\medskip
{\Large 採点基準}\par\medskip
2026年11月11日\quad 作成識別：6astra
\end{center}
\noindent 本冊子には解答の要点を含む。受験後に参照すること。
以下は本模試独自の配点であり，公式試験の採点基準を再現したものではない。

\section*{共通の扱い}
\begin{enumerate}[label=\arabic*.,leftmargin=2em,itemsep=0.7em]
\item 各大問30点。受験時に指定した3問の素点合計を $S$ とし，換算得点を $100S/90$ とする。
4問以上解いた場合は，採点を見る前に採点対象3問を決める。全5問の150点集計は復習用とする。
\item 各項目の点は加点式とし，小問の範囲内で採点する。数学的に正しい別解・同値な表現には同じ点を与える。
\item 前の小問の計算誤りを引き継いでも，その結果の下で後続の方法が正しければ，後続小問の方法に対応する点を与える。
同じ誤りを重複して機械的に減点しない。ただし誤った前提から本質的に別の問題になっている場合は，成立した部分だけを評価する。
\item 導出を要求した箇所で答えだけがある場合は，各欄の結論に対応する点までを与える。
表記の違いや自明な中間式の省略は減点しない。台・母数範囲の欠落は，結論の解釈や妥当性に関係する場合に対応する項目から減点する。
\item 数値は問題文に示した近似値を使用する。丸め前の正確な式が正しければ，末尾の軽微な丸め差は許容する。
問題に求められていない補足計算は加点・減点の対象としない。
\item 本模試の得点から公式試験の合否は判定しない。未選択問題は，採点済み3問とは別に復習する。
\end{enumerate}

\section*{配点一覧}
\begin{center}
\begin{tabular}{lrrrrrr}
\toprule
 &〔1〕&〔2〕&〔3〕&〔4〕&〔5〕&合計\\
\midrule
問1&4&7&5&8&6&30\\
問2&4&6&8&6&6&30\\
問3&8&9&8&5&---&30\\
問4&5&6&7&6&6&30\\
問5&5&6&6&6&7&30\\
\bottomrule
\end{tabular}
\end{center}
\newpage

\qhead{1}
\ans{1\ \textnormal{（4点）}}
条件付き母関数 $\E[s^W\mid N=j]=\{ps+(1-p)s^2\}^j$：2点。
ポアソン確率で平均し，$G_W(s)=\exp[\lambda\{ps+(1-p)s^2-1\}]$：2点。

\ans{2\ \textnormal{（7点）}}
値2を取る事象が $j$ 個なら $N=k-j$ と数える，または母関数を2つの指数関数の積に展開する方針：2点。
\[
\Prob(W=k)=e^{-\lambda}\sum_{j=0}^{\lfloor k/2\rfloor}
\frac{(\lambda p)^{k-2j}\{\lambda(1-p)\}^{j}}{(k-2j)!j!}
\]
を正しい和の範囲で示す：3点。
$\Prob(W=2)=e^{-\lambda}\{(\lambda p)^2/2+\lambda(1-p)\}$：2点。
確率母関数からポアソン分布の独立な2成分へ分解する解法も満点。

\ans{3\ \textnormal{（5点）}}
$\mu_W=\lambda(2-p)$ の導出：2点。
全分散の公式または母関数の微分で，$\sigma_W^2=\lambda(4-3p)$ の導出：3点。
条件付き分散の平均だけを計算して $\lambda p(1-p)$ とした場合，分散部分は最大1点。

\ans{4\ \textnormal{（8点）}}
$\hat\lambda=(3\bar W_n-V_n)/2$：1点。
$\hat p=2(2\bar W_n-V_n)/(3\bar W_n-V_n)$：2点。
条件 $\bar W_n>0,\ \bar W_n<V_n<2\bar W_n$ を必要十分として示す：3点（条件2点，両方向の確認1点）。
数値 $\hat\lambda=2,\hat p=1/2$：各1点。
分母が正という条件だけでは母数空間内に入る保証がないため，条件部分は最大1点。

\ans{5\ \textnormal{（6点）}}
極限 $N(0,\lambda(4-3p))$：2点。
独立同分布・有限な正の分散を確認して中心極限定理を適用：1点。
大数の法則で $\bar W_n,V_n$ の確率収束を示す：1点。
連続写像定理または同等の議論による両推定量の一致性：1点。
真値が条件の内部にあるため，条件外の標本の確率が0に近づくことを説明：1点。
\newpage

\qhead{2}
\ans{1\ \textnormal{（4点）}}
$X_i^2\sim\mathrm{Exp}(\theta)$（率）の変数変換と台：2点。
$\E[X_i^2]=1/\theta$：1点。
$B_i\sim\mathrm{Bernoulli}(e^{-\theta c^2})$：1点。

\ans{2\ \textnormal{（6点）}}
連続観測と二値観測の独立な寄与を正しく積にする：2点。
尤度の母数依存部分
\[
\theta^m e^{-\theta(s+kc^2)}(1-e^{-\theta c^2})^{r-k}
\]
を得る：2点。
因子分解を根拠として $(S,K)$ の十分性を示す：2点。
母数によらない正の因子を省略してもよい。二値観測を生の $X_i$ と同じ密度で扱うのは認めない。

\ans{3\ \textnormal{（8点）}}
正しい対数尤度：1点。
スコア方程式
\[
\frac m{\hat\theta}-(s+kc^2)+\frac{(r-k)c^2}{e^{\hat\theta c^2}-1}=0
\]
：2点。
$\ell''(\theta)<0$ を示す：2点。
$\theta\downarrow0$ でスコアが $+\infty$：1点。
$\theta\to\infty$ でスコアが $-(s+kc^2)<0$：1点。
以上から正の解の存在・一意性と大域的最大化を結論：1点。
「2階微分が負」のみでは存在を示していない。$k=r$ も含めて扱う。

\ans{4\ \textnormal{（6点）}}
$I_X=1/\theta^2$：1点。
$I_B=c^4/(e^{\theta c^2}-1)$：2点（スコア等の導出1点，結論1点）。
全情報量 $m/\theta^2+rc^4/(e^{\theta c^2}-1)$：1点。
比 $(\log2)^2$：1点。
問題文の値による $0.480$：1点。

\ans{5\ \textnormal{（6点）}}
最大化する関数 $h(z)=z^2/(e^z-1)$ とその導関数：2点。
方程式 $(2-z)e^z=2$：1点。
$g(z)=(2-z)e^z-2$ の増減等から，正の解が一意で $(1,2)$ にあり最大点となること：2点。
$c_*=\sqrt{z_*/\theta_0}$：1点。
$z=0$ は方程式の解だが許容範囲 $z>0$ に入らず，最適解ではない。
\newpage

\qhead{3}
\ans{1\ \textnormal{（8点）}}
$A,B$ の独立なガンマ密度から正しい同時密度を得る：2点。
$a=ut,b=(1-u)t$，$0<u<1,t>0$，ヤコビアン $t$：2点。
\[
f_U(u;\theta)=\frac{6\theta^2u(1-u)}{\{1+(\theta-1)u\}^4}
\]
：2点（積分1点，結果1点）。
$\theta=1$ の同時密度を，正規化された $U,T$ の密度の積に分解して独立性を示す：2点。
「共分散0」だけでは独立性の証明にならない。

\ans{2\ \textnormal{（9点）}}
元の観測モデルから尤度 $L=\lambda^4\theta^2e^{-\lambda(\theta a+b)}$ を得る：2点。
固定した $\theta$ で $\hat\lambda(\theta)=4/(\theta a+b)$ が最大化すること：2点。
プロファイル尤度の増減と制約から $\hat\theta=\max(1,b/a)$：3点（微分1点，増減1点，境界1点）。
\[
\Lambda=\begin{cases}\{4U(1-U)\}^2,&U<1/2,\\1,&U\geq1/2\end{cases}
\]
：2点。
$\theta\geq1$ を無視して両側用の比だけを出した場合，最後の2点は与えず，制約処理の1点も与えない。

\ans{3\ \textnormal{（8点）}}
帰無分布の累積確率 $3c^2-2c^3$：2点。
小さい $U$ 側の棄却域 $U\leq1/4$，サイズ $5/32$：2点。
一般の $\theta$ における検出力
\[
\pi(\theta)=\frac{\theta^2(\theta+9)}{(\theta+3)^3}
\]
：3点（変換または積分の方針1点，導出と結論2点）。
$\pi(3)=1/2$：1点。
連続分布なので $<$ と $\leq$ は同等。未知の $\lambda$ に依存しないことが式から明白なら説明の重複は不要。

\ans{4\ \textnormal{（5点）}}
$f_U(u;\theta_1)/f_U(u;1)=\theta_1^2/\{1+(\theta_1-1)u\}^4$ が $\theta_1>1$ で減少：2点。
ネイマン・ピアソンの基本定理を $U$ の分布に適用：1点。
臨界値が $\theta_1$ によらず同じなので一様最強力と結論：1点。
$U$ だけに基づく検定の中での主張として正しく述べる：1点。
「尤度比検定だから一様最強力」とだけ述べた場合は0点。
\newpage

\qhead{4}
\ans{1\ \textnormal{（5点）}}
KL情報量の積分による定義：2点。
正規密度の対数比と期待値の計算：2点。
$(\mu-\eta)^2/(2\sigma^2)$：1点。
この設問では分散が共通なのでKLの逆向きも値は同じだが，定義は指定された向きで採点する。

\ans{2\ \textnormal{（6点）}}
$R_n(a)=\{(1-a)^2\mu^2+a^2\sigma^2/n\}/(2\sigma^2)$：3点（偏り・分散の計算2点，結論1点）。
最小化と $a_{\mathrm{opt}}=n\mu^2/(n\mu^2+\sigma^2)$：2点。
未知の $\mu^2$ に依存するのでそのまま使えないこと：1点。
$\mu=0$ では $a_{\mathrm{opt}}=0$ となる式であれば，別記は必須でない。

\ans{3\ \textnormal{（7点）}}
$\mathrm{CV}_n(a)=n^{-1}\sum_i(Y_i-a\bar Y_{(-i)})^2$：2点。
$\bar Y_{(-i)}=(n\bar Y-Y_i)/(n-1)$：1点。
中心化した残差を使った変形，および交差項が消えること：2点。
\[
\mathrm{CV}_n(a)=\left(1+\frac a{n-1}\right)^2V_n+(1-a)^2\bar Y^2
\]
：2点。
全データ平均で評価した学習誤差をLOOCVとした場合，定義の2点は与えない。

\ans{4\ \textnormal{（6点）}}
\[
\E[\mathrm{CV}_n(a)]=\sigma^2+(1-a)^2\mu^2+a^2\sigma^2/(n-1)
\]
：3点（独立性等の根拠1点，式2点）。
$\E[(Y_0-a\bar Y)^2]=\sigma^2+(1-a)^2\mu^2+a^2\sigma^2/n$：2点。
差 $a^2\sigma^2/\{n(n-1)\}$：1点。
「予測誤差なので等しい」として学習標本数の違いを無視しない。

\ans{5\ \textnormal{（6点）}}
$\bar y=2,V_4=7/2$：1点。
$\mathrm{CV}_4(a)=15/2-(17/3)a+(79/18)a^2$ または同等の未展開式：1点。
制約を確認した最小点 $\hat a=51/79$：2点。
固定 $a$ の期待値公式を標本依存の $\hat a$ にそのまま代入できない理由：2点。
理由として，$\hat a$ と標本平均・分散の依存，評価用の $Y_i$ への依存，最小値と期待値の交換の問題等を正しく述べればよい。
「過学習するから」など具体的な依存関係の説明がないものは，理由部分は最大1点。
\newpage

\qhead{5}
\ans{1\ \textnormal{（5点）}}
$\E_*[\bar X^*]=\bar X$：1点。
条件付き独立性から $\Var_*(\bar X^*)=V_n/n$：2点。
$\E[\Var_*(\bar X^*)]=(n-1)\theta^2/(12n^2)$：1点。
$\Var(\bar X)=\theta^2/(12n)$ の $(n-1)/n$ 倍と比較：1点。
条件付きの分散と元の標本に関する無条件の分散を区別する。

\ans{2\ \textnormal{（6点）}}
$F_{M_n/\theta}(u)=u^n$（$0\leq u\leq1$，台の外の処理を含む）：2点。
$F_{T_n}(t)=1-(1-t/n)^n$（$0\leq t\leq n$，外では0または1）：2点。
$(1-t/n)^n\to e^{-t}$：1点。
率1の指数分布への法則収束：1点。

\ans{3\ \textnormal{（6点）}}
$\Prob_*(M_n^*\leq x)=F_n(x)^n$ と条件付き独立性：2点。
最大値の添え字を一度も選ばない事象を使う：1点。
$\Prob_*(M_n^*=M_n)=1-(1-1/n)^n$：2点。
極限 $1-e^{-1}$：1点。

\ans{4\ \textnormal{（6点）}}
$\Prob_*(T_n^*\leq0)=1-(1-1/n)^n$：1点。
真の極限分布関数は0で連続かつ値0なので収束に反するという論証：2点。
ブートストラップ分位点の上端も $M_n$ 以下：1点。
$M_n<\theta$ が確率1で成立：1点。
被覆確率0：1点。
「ブートストラップは離散だから失敗」とだけ述べるのは不十分。消えない原点の質量を使う。

\ans{5\ \textnormal{（7点）}}
$M_n^\dagger/M_n$ の条件付き分布関数が $u^n$：1点。
標準化後の分布関数が〔2〕と同じであることを台も含めて示す：2点。
被覆事象を $M_n/\theta\geq1/c_\alpha$ と表す：1点。
被覆確率 $1-c_\alpha^{-n}$：1点。
$c_\alpha=\alpha^{-1/n}$ と区間 $[M_n,M_n/\alpha^{1/n}]$：2点。
端点の開閉は連続分布のため同等。上端を $M_n\alpha^{1/n}$ とした区間は認めない。
\end{document}
